% Part: proof-theory
% Chapter: proof-search
% Section: tableaux

\documentclass[../../../include/open-logic-section]{subfiles}

\begin{document}

\olfileid[id]{pt}{ps}{tab}

\olsection{Tableau}

Sistem tableau adalah sistem !!{proof} yang berkaitan erat dengan
kalkulus sekuen dan sangat sesuai untuk pencarian bukti, baik dengan
tangan maupun melalui implementasi perangkat lunak. Alih-alih membangun
pohon sekuen, dalam tableau !!a{proof} merupakan pohon !!{formula}s.
Namun, tidak seperti deduksi alami, aturan-aturannya mencerminkan aturan
kalkulus sekuen. !!{proof} tableau pada dasarnya merupakan pencarian
eksplisit atas suatu model. Karena itu, sistem ini kadang-kadang disebut
tableau \emph{semantik} atau juga \emph{pohon kebenaran}.

Tableau \emph{bertanda} adalah pohon !!{formula}s bertanda, yang dapat
berbentuk $\sFmla{\True}{!A}$ atau $\sFmla{\False}{!A}$. Pohon tumbuh ke
bawah. Pohon dimulai dengan !!{formula}s yang menentukan apa yang ingin
kita !!{prove}. Misalnya, untuk !!{prove} sekuen $!A, !B \Sequent !C,
!D$, tableau dimulai dengan $\sFmla{\True}{!A}, \sFmla{\True}{!B},
\sFmla{\False}{!C}, \sFmla{\False}{!D}$. !!^a{structure} yang membuat
$!A$ dan $!B$ benar serta $!C$ dan $!D$ salah adalah !!a{structure} yang
menunjukkan bahwa $!A, !B \Sequent !C, !D$ tidak sahih.

Simpul daun (paling bawah) suatu tableau dapat diperluas dengan aturan
perluasan cabang, yang menambahkan simpul lain pada cabang yang sama atau
menambahkan dua simpul bersaudara di bawahnya sehingga cabang terbelah
menjadi dua cabang baru. Aturan dapat diterapkan kepada sembarang
!!{formula} bertanda di atas simpul itu pada cabang yang sama. Aturan
perluasan cabang bagi sistem tableau klasik \Log{Tc} diberikan dalam
\olref{tab:Tc}.

\subfile{rules-Tc}

Seperti dapat dilihat, aturan-aturan tersebut hanyalah aturan \Log{G3c}
yang dibalik, dan tanda $\True$ serta $\False$ masing-masing menunjukkan
apakah !!a{formula} menjadi formula utama di sisi kiri atau kanan suatu
sekuen.

Sebuah cabang tableau \emph{tertutup} jika memuat baik
$\sFmla{\True}{!A}$ maupun $\sFmla{\False}{!A}$. Jika setiap cabang
tertutup, seluruh tableau disebut tertutup. Misalnya, berikut tableau
klausul untuk $!A \lor (!B \land !C) \Sequent (!A \lor !B) \land
(!A \lor !C)$, yakni dimulai dengan $\sFmla{\True}{!A \lor (!B \land
!C)}$ dan $\sFmla{\False}{(!A \lor !B) \land (!A \lor !C)}$:
\begin{oltableau}
  [\sFmla{\True}{\formula{A} \lor (\formula{B} \land \formula{C})},
    just=\TAss, checked
    [\sFmla{\False}{(\formula{A} \lor \formula{B}) \land
        (\formula{A} \lor \formula{C})}, just=\TAss, checked
      [\sFmla{\True}{\formula{A}}, just={\TRule{\True}{\lor}[1]}
        [\sFmla{\False}{\formula{A} \lor \formula{B}},
          just={\TRule{\False}{\land}[2]}, checked
          [\sFmla{\False}{\formula{A}},
            just={\TRule{\False}{\lor}[4]}
            [\sFmla{\False}{\formula{B}},
              just={\TRule{\False}{\lor}[4]}, close]
          ]
        ]
        [\sFmla{\False}{\formula{A} \lor \formula{C}},
          just={\TRule{\False}{\land}[2]}, checked
          [\sFmla{\False}{\formula{A}},
            just={\TRule{\False}{\lor}[4]}
            [\sFmla{\False}{\formula{C}},
              just={\TRule{\False}{\lor}[4]}, close]
          ]
        ]
      ]
      [\sFmla{\True}{\formula{B} \land \formula{C}},
        just={\TRule{\True}{\lor}[1]}, checked
        [\sFmla{\False}{\formula{A} \lor \formula{B}},
          just={\TRule{\False}{\land}[2]}, checked
          [\sFmla{\True}{\formula{B}},
            just={\TRule{\True}{\land}[3]}, move by=2
            [\sFmla{\True}{\formula{C}},
              just={\TRule{\True}{\land}[3]}
              [\sFmla{\False}{\formula{A}},
                just={\TRule{\False}{\lor}[4]}
                [\sFmla{\False}{\formula{B}},
                  just={\TRule{\False}{\lor}[4]},close
                ]
              ]
            ]
          ]
        ]
        [\sFmla{\False}{\formula{A} \lor \formula{C}},
          just={\TRule{\False}{\land}[2]}, checked
          [\sFmla{\True}{\formula{B}},
            just={\TRule{\True}{\land}[3]}, move by=2
              [\sFmla{\True}{\formula{C}},
                just={\TRule{\True}{\land}[3]}
                [\sFmla{\False}{\formula{A}},
                  just={\TRule{\False}{\lor}[4]}
                  [\sFmla{\False}{\formula{C}},
                  just={\TRule{\False}{\lor}[4]},close
                ]
              ]
            ]
          ]
        ]
      ]
    ]
  ]
\end{oltableau}

Tanda centang menunjukkan bahwa aturan yang bersangkutan telah diterapkan
kepada !!a{formula} pada semua cabang di bawahnya. Simbol $\otimes$
menunjukkan bahwa suatu cabang tertutup. Nomor baris setelah aturan
menunjukkan kepada !!{formula} bertanda mana aturan itu diterapkan (yakni
!!{formula} pada cabang yang sama di baris yang ditunjukkan).

Pencarian bukti mudah dilakukan dalam tableau. Kita menghasilkan tableau
secara bertahap. Pada tahap~$0$, tuliskan saja !!{formula}s awal di bagian
atas. Pada tahap~$n+1$, lakukan hal berikut untuk setiap simpul daun:
Carilah formula bertanda pertama (paling atas) yang belum dicentang.
Terapkan aturan perluasan cabang kepada formula tersebut. Setelah aturan
diterapkan pada setiap cabang yang memuat formula bertanda itu, beri tanda
centang. (Contoh di atas dihasilkan dengan algoritma ini.)

!!{formula}s bertanda yang berbentuk $\sFmla{\True}{\lforall}$ dan
$\sFmla{\False}{\lexists}$ tidak pernah dicentang, sehingga harus
diperlakukan secara khusus. Jika formula-formula itu ada, kita harus
memastikan bahwa aturan $\TRule{\True}{\forall}$ dan
$\TRule{\False}{\lexists}$ akhirnya diterapkan untuk semua suku yang
berlaku. Hal ini dapat dijamin, misalnya, dengan menerapkannya pada setiap
tahap kepada semua formula semacam itu di setiap cabang (setelah menangani
!!{formula} bertanda paling atas yang tidak berbentuk demikian). Suku $t$
yang digunakan ialah suku pertama~$t$ yang dapat dibangun dari !!{constant}s
yang sudah muncul pada cabang (atau $\Obj c_0$ jika tidak ada
!!{constant} yang muncul) dan !!{function}s yang muncul dalam
!!{formula}s awal, sedemikian rupa sehingga !!{formula} bertanda yang
bersesuaian, $\sFmla{\True}{!B(t)}$ atau~$\sFmla{\False}{!B(t)}$, belum
muncul pada cabang.

Jika pencarian bukti ini tidak menghasilkan tableau tertutup, tableau itu
memuat cabang terbuka berhingga tempat semua !!{formula} nonatomiknya telah
dicentang, atau pencarian menghasilkan pohon takhingga bercabang hingga,
yang harus memuat suatu cabang takhingga. Sama seperti dalam
\olref[cpl]{sec}, kita dapat mendefinisikan !!a{structure}~$\Struct{M}$
sedemikian rupa sehingga jika $\sFmla{\True}{!A}$ muncul pada cabang,
maka $\Sat{M}{!A}$, dan jika $\sFmla{\False}{!A}$ muncul, maka
$\Sat/{M}{!A}$. (Ambil $\Theta =
\Setabs{!A}{\sFmla{\True}{!A} \text{ pada cabang}}$ dan $\Xi =
\Setabs{!A}{\sFmla{\False}{!A} \text{ pada cabang}}$.)

Tableau yang dihasilkan oleh metode pencarian bukti ini dapat dengan mudah
diterjemahkan menjadi bukti \Log{G3c} dengan memberikan sebuah sekuen
kepada setiap simpul. Jika formula-formula awal bersesuaian dengan sekuen
$\Gamma \Sequent \Delta$, beri setiap formula awal label $\Gamma \Sequent
\Delta$. Jika suatu cabang diperluas dengan menerapkan aturan perluasan
cabang kepada !!{formula} bertanda $\sFmla{\True}{!A}$, daun itu diberi
label $!A, \Pi \Sequent \Lambda$; jika aturan diterapkan kepada
!!{formula} bertanda $\sFmla{\False}{!A}$, daun itu diberi label
$\Pi \Sequent \Lambda, !A$. Ketika tableau menerapkan aturan
\TRule{\True}{\star}, terapkan aturan \LeftR{\star} kepada sekuen (dengan
$!A$ sebagai formula utama); ketika tableau menerapkan aturan
\TRule{\False}{\star}, terapkan aturan \RightR{\star}. Premis-premis aturan
tersebut ialah sekuen yang melabeli simpul-simpul bersesuaian yang
ditambahkan kepada tableau. Ketika suatu tableau tertutup oleh
$\sFmla{\True}{!A}$ dan $\sFmla{\False}{!A}$, label simpul daun terakhir
akan berupa sekuen $!A, \Pi \Sequent \Lambda, !A$. Beberapa simpul
berturutan dalam tableau akan diberi sekuen yang sama; hapus duplikat.
Misalnya, tableau di atas diterjemahkan menjadi !!{proof} berikut:
\[\small
\Axiom$!A \fCenter !A, !B$
\RightLabel{\RightR{\lor}}
\UnaryInf$!A \fCenter !A \lor !B$
\Axiom$!A \fCenter !A, !C$
\RightLabel{\RightR{\lor}}
\UnaryInf$!A \fCenter !A \lor !C$
\RightLabel{\RightR{\land}}
\BinaryInf$!A \fCenter (!A \lor !B) \land (!A \lor !C)$
\Axiom$!B, !C \fCenter !A, !B$
\RightLabel{\RightR{\lor}}
\UnaryInf$!B, !C \fCenter !A \lor !B$
\RightLabel{\LeftR{\land}}
\UnaryInf$!B \land !C \fCenter !A \lor !B$
\Axiom$!B, !C \fCenter !A, !C$
\RightLabel{\RightR{\lor}}
\UnaryInf$!B, !C \fCenter !A \lor !C$
\RightLabel{\LeftR{\land}}
\UnaryInf$!B \land !C \fCenter !A \lor !C$
\RightLabel{\RightR{\land}}
\BinaryInf$!B \land !C \fCenter (!A \lor !B) \land
(!A \lor !C)$
\RightLabel{\LeftR{\lor}}
\BinaryInf$!A \lor (!B \land !C) \fCenter (!A \lor !B) \land
(!A \lor !C)$
\DisplayProof
\]

\end{document}
